% 九秒教学序列，样本线不表示实际方法的固定帧数。
\begin{tikzpicture}[x=0.82cm,y=1cm,font=\figurefont\small]
  \foreach \a/\b/\name in {0/3/伸手,3/6/拿杯,6/9/收手}{
    \fill[BookPaper] (\a,2.45) rectangle (\b,3.05);
    \node at ({(\a+\b)/2},2.75) {\name};
    \draw[BookRule,line width=0.5pt] (\a,0.2) -- (\a,3.1);
  }
  \draw[BookRule,line width=0.5pt] (9,0.2) -- (9,3.1);
  \draw[-{Stealth[length=2mm]},BookInk,line width=0.8pt]
    (0,0) -- (9.3,0);
  \node[anchor=east] at (9.3,-0.65) {时间（秒）};
  \foreach \x in {0,3,6,9}{
    \draw[BookInk,line width=0.7pt] (\x,0.05) -- (\x,-0.05);
    \node[below] at (\x,-0.08) {\x};
  }
  \node[anchor=east] at (-0.2,1.95) {跨段取样};
  \draw[BookMuted,line width=0.6pt] (0,1.95) -- (9,1.95);
  \foreach \x in {1.5,4.5,7.5}{
    \draw[BookBlue,line width=1.5pt] (\x,1.7) -- (\x,2.2);
  }
  \node[anchor=east] at (-0.2,0.95) {连续片段};
  \draw[BookMuted,line width=0.6pt] (0,0.95) -- (9,0.95);
  \fill[BookTeal!12] (3.6,0.65) rectangle (5.4,1.25);
  \foreach \x in {3.7,3.95,4.2,4.45,4.7,4.95,5.2}{
    \draw[BookTeal,line width=1pt] (\x,0.7) -- (\x,1.2);
  }
\end{tikzpicture}
